% C15 — Qanat underground watercourse (I-8). Book: vol-1/ch-08 ¶ "Come as far as the end of that third day in Ramusio's account", Note [3] ([@yc, 1.20 n. 2]).
% Pure TikZ cross-section per DIA.C15 specification.
\path (0,0) -- (11.43,0);
\node[dg-head] at (0,0.25) {Underground watercourse (\textit{qanāt} / \textit{kārīz}), as Yule explains the desert stream};
% Ground profile and subsurface fill
\fill[inksoft!15] (0.2,1.95) -- (2.2,2.45) -- (8.8,3.45) -- (11.2,3.65) -- (11.2,4.65) -- (0.2,4.65) -- cycle;
\draw[ink, line width=0.6pt] (0.2,1.95) -- (2.2,2.45) -- (8.8,3.45) -- (11.2,3.65);
% Water table at piedmont
\fill[vIII!20] (0.2,2.95) .. controls (1.2,2.95) and (2.0,3.15) .. (2.6,3.45) -- (0.2,3.45) -- cycle;
\draw[vIII, line width=0.5pt, dashed] (0.2,2.95) .. controls (1.2,2.95) and (2.0,3.15) .. (2.6,3.45);
% Gently sloping underground gallery
\fill[white] (1.3,3.25) -- (9.3,3.70) -- (9.3,3.92) -- (1.3,3.47) -- cycle;
\fill[vIII!65] (1.3,3.37) -- (9.8,3.85) -- (9.8,3.94) -- (1.3,3.47) -- cycle;
\draw[ink, line width=0.5pt] (1.3,3.25) -- (9.3,3.70);
\draw[ink, line width=0.5pt] (1.3,3.47) -- (9.8,3.94);
% Vertical shafts with spoil rings at surface
\foreach \x/\yg/\yb in {1.35/2.24/3.26, 3.10/2.59/3.36, 4.85/2.85/3.46, 6.60/3.12/3.56, 8.35/3.38/3.66} {
  \fill[white] ({\x-0.09},\yg) rectangle ({\x+0.09},\yb);
  \draw[inksoft, line width=0.4pt] ({\x-0.09},\yg) -- ({\x-0.09},\yb);
  \draw[inksoft, line width=0.4pt] ({\x+0.09},\yg) -- ({\x+0.09},\yb);
  \fill[ink] ({\x-0.18},\yg) -- ({\x-0.09},{\yg-0.10}) -- ({\x-0.09},\yg) -- cycle;
  \fill[ink] ({\x+0.18},\yg) -- ({\x+0.09},{\yg-0.10}) -- ({\x+0.09},\yg) -- cycle;
}
% Labels
\draw[tl-leader] (1.35,2.10) -- (1.35,1.55);
\node[dg-text, text width=3.4cm, font=\scriptsize, anchor=north west] at (0.2,0.75) {\textbf{Mother well} taps the alluvial water table at the hill foot};
\draw[tl-leader] (4.85,2.70) -- (5.35,1.75);
\node[dg-text, text width=4.8cm, font=\scriptsize, anchor=north west] at (4.2,0.75) {\textbf{Vertical shafts}: Ramusio's text notes openings where travellers see the buried stream and draw water};
\draw[tl-leader] (9.55,3.85) -- (9.55,2.65);
\node[dg-text, text width=2.6cm, font=\scriptsize, anchor=north west] at (8.6,1.85) {\textbf{Outlet} to surface channel and fields};
\node[dg-note, text width=11.2cm] at (0,4.85) {In Ramusio's account of the Kerman–Cobinan desert, a fresh stream runs underground and is reached at holes along its course. Yule identifies it as a Persian \textit{kārīz} (\textit{qanāt}).};
